% =============================================================================
% Appendix I  Verification Protocol  --  cited from sec:pipeline-reference and
% sec:pipeline-verification.   Plan: ThesisPlanning.md §6 (I).
%
% CONTENT  The verification modes, sizes and tolerances; the structured probes;
%          per-atom residuals; the circuit-simulation memory and time model and
%          its calibration; the acceptance criteria and the error-budget
%          measurement; the eigh-against-SVD comparison.
% SOURCES  Planning.md §6.2, §7.3, §7.8-7.9, App. A-3; Journal E023-E024;
%          pihm/src/pihm/verify.py, classical.py, qsp/prepare.py.
% =============================================================================

\subsection{Modes, sizes and tolerances}
\label{app:verification-modes}
% LAND: tab:verification-modes in full -- how exact mode extracts
%   Pi_out U Pi_in column by column, how probe mode estimates ||E||_F from Haar
%   and structured inputs, what IR mode compares at every n, and why it is not
%   a tautology (the operator it checks against is assembled independently).

\subsection{Structured probes}
\label{app:structured-probes}
% LAND: the probes that target where past bugs lived -- cyclic wrap-around,
%   T_0 weights, block edges, axis order, borrow flags (lessons L-01 to L-03,
%   L-20) -- stated as coverage, not as history.

\subsection{Per-atom residuals}
\label{app:atom-residuals}
% LAND: a table of verified residuals per atom and n. Include the banded-
%   factorisation residual ||Bt^{-1} Dt - G|| over n = 1..10 from pihm's R0
%   property tests (the draft's TODO named the legacy test
%   test_descriptor_encoding.py::test_banded_factorisation_is_exact; cite the
%   rebuilt package's test instead).

\subsection{The classical reference: eigh against SVD}
\label{app:eigh-vs-svd}
% The values below are the planning probe's (Planning.md App. A-3: {1,4,4} with
%   f(-1/2) = 0, reference (1/2 + x) e^{-2x}); regenerate at results-v1.

\begin{table}[htbp]
    \centering
    \caption{Accuracy of the classical reference, $1 - |\cos|$ against the exact coefficients, for the standard form by \texttt{eigh}$(R^{\top}R)$ and by SVD of $R$, and for the descriptor form.}
    \label{tab:eigh-vs-svd}
    \small
    \begin{tabular}{@{}r l l l@{}}
        \toprule
        $n$ & Standard, \texttt{eigh} & Standard, SVD & Descriptor \\
        \midrule
        2 & $\prov{1.7\times10^{-1}}$  & $\prov{1.7\times10^{-1}}$ & $\prov{1.6\times10^{-2}}$ \\
        3 & $\prov{8.6\times10^{-8}}$  & $\prov{8.6\times10^{-8}}$ & $\prov{5.8\times10^{-9}}$ \\
        4 & $\prov{0}$                  & $\prov{\le 2\times10^{-16}}$ & $\prov{\le 2\times10^{-16}}$ \\
        5 & $\prov{3.7\times10^{-15}}$ & $\prov{\le 2\times10^{-16}}$ & $\prov{\le 2\times10^{-16}}$ \\
        6 & $\prov{1.7\times10^{-10}}$ & $\prov{\le 2\times10^{-16}}$ & $\prov{\le 2\times10^{-16}}$ \\
        7 & $\prov{1.2\times10^{-4}}$  & $\prov{\le 2\times10^{-16}}$ & $\prov{\le 2\times10^{-16}}$ \\
        8 & $\prov{0.73}$               & $\prov{\le 2\times10^{-16}}$ & $\prov{\le 2\times10^{-16}}$ \\
        \bottomrule
    \end{tabular}
\end{table}

\subsection{The circuit-simulation memory and time model}
\label{app:t1-model}
% LAND: peak memory = the statevector, 16 * 2^q bytes for q qubits, plus one
%   walk's own gates, estimated from the circuit's structure; time = d x (one
%   measured step) x (initial-state candidates); `infeasible` recorded with the
%   estimate and T3's numbers when either exceeds the run's budget (decision
%   20). Calibration: measured peak / model = 0.85-1.05 at 22 qubits (Journal
%   E024). The laptop benchmarks: R1 at n = 2 (11 qubits, d = 89) in 0.3 s;
%   R2a at n = 2 (20 qubits, d = 544) at 4.2 s per step (\prov).

\subsection{Acceptance and the error budget}
\label{app:error-budget}
% LAND: tab:acceptance applied run by run; each term of eq:error-budget measured
%   separately, and how (eps_disc from the classical reference against the
%   analytic solution; eps_filter from the ideal filter; eps_trunc and
%   eps_phase from the realised polynomial; d eps_BE for the approximate
%   published G only; eps_sample from the shot count).
